\documentclass[
  aps,
  pre,
  reprint,
  superscriptaddress,
  nofootinbib
]{revtex4-2}

\usepackage[T1]{fontenc}
\usepackage{lmodern}
\usepackage{amsmath,amssymb,amsfonts,mathtools}
\usepackage{booktabs}
\usepackage{hyperref}

\hypersetup{
  colorlinks=true,
  linkcolor=blue,
  citecolor=blue,
  urlcolor=blue
}

\newcommand{\diag}{\operatorname{diag}}
\newcommand{\Half}{\mathrm{half}}

\begin{document}

\title{Summary of the Non-Diagonal-Mass Construction Path}
\author{SphericalSBPOperators Development Team}
\affiliation{Computational Science and Engineering}
\date{\today}

\begin{abstract}
This note summarizes the current \texttt{src/non\_diagonal\_mass/} construction path in
\texttt{SphericalSBPOperators.jl}.
The emphasis is on the ansatz for the scalar mass $S$ and vector mass $V$, and on the
conditions used in \texttt{sbp4\_solve\_accuracy\_constraints},
\texttt{sbp6\_solve\_accuracy\_constraints}, and
\texttt{sbp8\_solve\_accuracy\_constraints} to determine them.
\end{abstract}

\maketitle

\section{Common construction path}
All three non-diagonal-mass paths follow the same top-level sequence:
\begin{enumerate}
  \item Build a full-grid Cartesian SBP operator on $[-R,R]$.
  \item Fold it to the half-grid $r \in [0,R]$ using even/odd parity to obtain
  \begin{equation}
    G_{\mathrm{even}} = R_{\mathrm{op}} G_{\mathrm{full}} E_{\mathrm{even}},
    \qquad
    G_{\mathrm{odd}} = R_{\mathrm{op}} G_{\mathrm{full}} E_{\mathrm{odd}}.
  \end{equation}
  \item Form the folded Cartesian half-mass
  \begin{equation}
    H_{\mathrm{cart},\Half}
    = \frac12 E_{\mathrm{even}}^{T} H_{\mathrm{full}} E_{\mathrm{even}}.
  \end{equation}
  \item Use the metric weight $r^p$ to define the diagonal seed
  \begin{equation}
    S_{\mathrm{seed}}
    = H_{\mathrm{cart},\Half}\,\diag(r^p).
  \end{equation}
  \item Solve for a diagonal scalar mass $S$ and a structured symmetric vector mass $V$,
  then define the divergence through the SBP identity
  \begin{equation}
    D = S^{-1}(B - G_{\mathrm{even}}^{T} V),
    \qquad
    B = \diag(0,\dots,0,r_N^p).
  \end{equation}
\end{enumerate}

The public constructors
\texttt{sbp4\_operators}, \texttt{sbp6\_operators}, and \texttt{sbp8\_operators}
all return
\begin{equation}
  (D,G,S,V,B), \qquad G = G_{\mathrm{even}}.
\end{equation}

\section{What is assumed for \texorpdfstring{$S$}{S} and \texorpdfstring{$V$}{V}?}
\subsection{Assumptions shared by all orders}
\begin{itemize}
  \item $S$ is always diagonal.
  \item $V$ is always symmetric.
  \item The non-diagonal part of $V$ is localized near the origin.
  \item Away from the origin, the code ties the tail of $V$ back to $S$ through
  $V_{ii}=S_{ii}$ once the boundary block ends.
\end{itemize}

So the entire method can be read as:
\begin{quote}
Keep $S$ simple, keep the far field diagonal, and put the non-diagonal freedom only in a
small origin block of $V$.
\end{quote}

\subsection{SBP4 ansatz}
The symbolic template in \texttt{sbp4\_vector\_mass} assumes
\begin{equation}
  S = \diag(s_1,\dots,s_N).
\end{equation}
For $V$:
\begin{align}
  V_{11} &= 1, \\
  V_{22},V_{33},V_{44},V_{55} &\text{ are free}, \\
  V_{23}=V_{32},\;
  V_{34}=V_{43}, \\
  V_{45}=V_{54} &\text{ is the last allowed off-diagonal pair}, \\
  V_{ii} &= s_i \quad \text{for } i \ge 6.
\end{align}
Thus only the leading $5 \times 5$ block can differ from the diagonal seed.

\subsection{SBP6 ansatz}
The symbolic template in \texttt{sbp6\_vector\_mass} keeps
\begin{equation}
  S = \diag(s_1,\dots,s_N),
\end{equation}
but expands the admissible boundary block for $V$:
\begin{align}
  V_{11} &= 1, \\
  V_{ii} &\text{ are free for } i=2,\dots,\min(10,N), \\
  V_{ii} &= s_i \quad \text{for } i \ge 11.
\end{align}
The allowed off-diagonal pairs are
\begin{align}
  &(2,3),(2,4), \nonumber\\
  &(3,4),(3,5), \nonumber\\
  &(4,5),(4,6), \nonumber\\
  &(5,6),(5,7), \nonumber\\
  &(6,7),
\end{align}
whenever those indices exist.
So the SBP6 $V$ block is diagonal in the tail and pentadiagonal-like near the origin.

\subsection{SBP8 ansatz}
For SBP8 the structure is more flexible.
The template in \texttt{sbp8\_vector\_mass} assumes
\begin{equation}
  S = \diag(s_1,\dots,s_N),
\end{equation}
and lets $V$ have a banded symmetric boundary block:
\begin{align}
  V_{ii} &\text{ free for } i \le n_b, \\
  V_{ii} &= s_i \quad \text{for } i > n_b,
\end{align}
with off-diagonals allowed only inside the leading
$n_b \times n_b$ block and only up to a chosen odd bandwidth.
The default exact path uses a heptadiagonal pattern
(\texttt{boundary\_bandwidth = 7}).

Unlike SBP4 and SBP6, SBP8 does not hard-wire $V_{11}=1$.
The origin diagonal of $V$ is solved together with the rest of the boundary block.

\section{What conditions are used to solve for \texorpdfstring{$S$}{S} and \texorpdfstring{$V$}{V}?}
\subsection{Shared algebraic condition}
All three orders use the same discrete divergence identity:
\begin{equation}
  S D + G_{\mathrm{even}}^{T} V = B.
\end{equation}
Applied to an odd monomial $u(r)=r^d$, this gives the row-wise condition
\begin{equation}
  -\bigl(G_{\mathrm{even}}^{T} V u\bigr)_i
  - s_i (p+d) r_i^{d-1}
  + B_{ii} u_i = 0.
  \label{eq:div-row}
\end{equation}
This is linear in the unknown entries of $S$ and $V$, and it is the core equation used
to populate the solve matrix in every order.

All three orders also impose the scalar quadrature condition
\begin{equation}
  \sum_{i=1}^{N} s_i = \int_0^R r^p\,dr = \frac{R^{p+1}}{p+1}.
  \label{eq:quad}
\end{equation}

\subsection{SBP4 conditions}
The exact SBP4 solve only frees the boundary block.
The unknown vector is
\begin{equation}
  x =
  [s_1,\dots,s_5,\,
   v_2,v_3,v_4,v_5,\,
   v_{23},v_{34},v_{45}]^{T},
\end{equation}
while for $i \ge 6$ the scalar tail is fixed to
\begin{equation}
  s_i = \bigl(S_{\mathrm{seed}}\bigr)_{ii},
  \qquad
  V_{ii}=s_i.
\end{equation}

The solver then enforces:
\begin{align}
  D r &= (p+1) \quad \text{on rows } 1:(N-4), \\
  D r^3 &= (p+3) r^2 \quad \text{on rows } 1:\min(5,N), \\
  \sum_i s_i &= \frac{R^{p+1}}{p+1}.
\end{align}

So SBP4 assumes that only the first five scalar masses and the first few vector-mass
entries need to move; the far field is inherited directly from the folded diagonal seed.

\subsection{SBP6 conditions}
The SBP6 exact solve frees the entire diagonal of $S$:
\begin{equation}
  x =
  [s_1,\dots,s_N,\,
   v_2,\dots,v_{\min(10,N)},\,
   \{v_{ij}\}_{(i,j)\in\mathcal{P}_6}]^{T},
\end{equation}
where $\mathcal{P}_6$ is the allowed off-diagonal pair list above.
For $i \ge 11$, the vector-mass diagonal is tied to the solved scalar mass:
\begin{equation}
  V_{ii} = s_i.
\end{equation}

The row sets are chosen from the inferred right-boundary closure width
\texttt{closure\_right}, and the solver enforces:
\begin{align}
  D r &= (p+1)
  && \text{on all rows}, \\
  D r^3 &= (p+3) r^2
  && \text{on rows } 1:(N-\texttt{closure\_right}), \\
  D r^5 &= (p+5) r^4
  && \text{on rows } 1:\min(\texttt{first\_rows\_r5},N), \\
  \sum_i s_i &= \frac{R^{p+1}}{p+1}.
\end{align}

Relative to SBP4, SBP6 relaxes much more of the diagonal structure.
The code is effectively assuming that sixth order needs a wider and more global
redistribution of scalar weights.

\subsection{SBP8 exact conditions}
The preferred SBP8 path is the exact rational solve in
\texttt{\_sbp8\_solve\_accuracy\_constraints\_exact}.
It frees:
\begin{itemize}
  \item all scalar masses $s_1,\dots,s_N$,
  \item the selected leading vector diagonal entries
  $V_{ii}$ for $i \in \{1,\dots,v_{\mathrm{diag,end}}\}$,
  \item the allowed boundary-block off-diagonals inside the banded pattern.
\end{itemize}
Outside that free diagonal set, the code ties $V_{ii}$ back to $s_i$.

The enforced polynomial conditions are
\begin{align}
  D r &= (p+1)
  && \text{on all rows}, \\
  D r^3 &= (p+3) r^2
  && \text{on rows excluding the right closure}, \\
  D r^5 &= (p+5) r^4
  && \text{on rows excluding the right closure}, \\
  D r^7 &= (p+7) r^6
  && \text{on rows } 1:\min(\texttt{first\_rows\_r7},N), \\
  \sum_i s_i &= \frac{R^{p+1}}{p+1}.
\end{align}

The exact SBP8 solve also adds an origin anchor
\begin{equation}
  s_1 = s_{1,\mathrm{target}}
\end{equation}
to avoid degenerate exact solutions with $S_{11}=0$.
By default this target is chosen from the local grid spacing,
approximately $0.55\,h^3$ when $p=2$.

\section{Extra admissibility conditions in the SBP8 path}
SBP8 is the only order that has a second solve path if the exact solve fails.
The default public function tries the exact solve first and falls back to a
JuMP/Ipopt nonlinear program.

In the fallback path, the equality constraints are the same divergence and quadrature
conditions as above, but the code additionally enforces:
\begin{enumerate}
  \item positivity bounds
  \begin{equation}
    s_i \ge s_{\min}, \qquad V_{ii} \ge s_{\min}
  \end{equation}
  on the free boundary variables;
  \item positive definiteness of the boundary block of $V$ through a Cholesky
  factorization
  \begin{equation}
    V_{\mathrm{boundary}} = L L^T;
  \end{equation}
  \item a small objective that keeps $S$ near the folded seed and keeps $V$ close to $S$.
\end{enumerate}

After either solve path, SBP8 also checks:
\begin{itemize}
  \item positivity of the diagonal of $S$,
  \item positive definiteness of the boundary block of $V$,
  \item optionally positive definiteness of the full $V$,
  \item floating-point level satisfaction of the enforced constraints,
  \item optionally that the spectrum of $L = D G$ is real and non-positive to within
  tolerance.
\end{itemize}

\section{Takeaway}
The non-diagonal-mass path can be summarized in one sentence:
\begin{quote}
Build a folded half-grid gradient, keep the scalar mass diagonal, allow only a compact
origin-local non-diagonal correction in the vector mass, and determine the free
coefficients by enforcing odd-monomial divergence exactness plus one scalar quadrature
constraint.
\end{quote}

Order by order, the main difference is how much freedom is granted:
\begin{itemize}
  \item SBP4: only the first five scalar masses and a tiny $V$ boundary block move.
  \item SBP6: all scalar masses move, with a wider pentadiagonal-like $V$ boundary block.
  \item SBP8: all scalar masses move, the $V$ boundary block becomes general banded, and
  the code adds an origin anchor plus stronger admissibility checks.
\end{itemize}

\end{document}
